\documentclass[11pt,a4paper]{article}

% ---------- Idioma y codificación ----------
\usepackage[spanish,es-tabla]{babel}
\usepackage[utf8]{inputenc}
\usepackage[T1]{fontenc}
\usepackage{lmodern}

% ---------- Formato de página ----------
\usepackage[margin=2.5cm]{geometry}
\usepackage{setspace}
\setstretch{1.15}
\usepackage{parskip}

% ---------- Tablas y listas ----------
\usepackage{booktabs}
\usepackage{tabularx}
\usepackage{array}
\usepackage{enumitem}

% ---------- Títulos en negrita (sin color) ----------
\usepackage{titlesec}
\titleformat{\section}
  {\Large\bfseries}
  {\thesection}{0.8em}{}
\titleformat{\subsection}
  {\large\bfseries}
  {\thesubsection}{0.6em}{}
\titleformat{\subsubsection}
  {\normalsize\bfseries}
  {\thesubsubsection}{0.5em}{}
\titlespacing*{\section}{0pt}{18pt}{8pt}

% ---------- Encabezado y pie ----------
\usepackage{fancyhdr}
\pagestyle{fancy}
\fancyhf{}
\fancyhead[L]{\small\textbf{Scrum --- Proyecto completo}}
\fancyhead[R]{\small\textbf{CrediStep}}
\fancyfoot[C]{\small\thepage}
\renewcommand{\headrulewidth}{0.4pt}

% ---------- Caja del objetivo (solo blanco y negro) ----------
\usepackage[most]{tcolorbox}
\newtcolorbox{objetivo}{
  colback=white, colframe=black,
  boxrule=0.8pt, arc=0pt,
  left=8pt, right=8pt, top=6pt, bottom=6pt,
  fonttitle=\bfseries, title={Objetivo del Proyecto}
}

% ======================================================
\begin{document}

% ---------- Portada ----------
\begin{center}
  {\Huge\bfseries Documento de Proyecto} \\[0.4cm]
  {\LARGE\bfseries Sprints 1 a 10} \\[0.3cm]
  {\large CrediStep} \\[1.2cm]
\end{center}

% ---------- 1. Información general ----------
\section{Información general}

\begin{tabularx}{\textwidth}{>{\bfseries}l X}
  \toprule
  Proyecto             & CrediStep \\
  Asignatura           & Gestión de Innovación en TI, Pontificia Universidad Javeriana \\
  Product Owner        & Víctor Hugo Vela Pinzón (profesor) \\
  Scrum Master         & Beiker Romario Jaramillo Giraldo \\
  Equipo de desarrollo & Laura Sofía Ortiz Gómez, Sarahy Michelle Silva Pinto, Valentina Carrillo Peñuela, Miguel Ángel Gómez López \\
  Fecha de inicio      & 16/07/2026 \\
  Fecha de fin         & 25/11/2026 \\
  Número de sprints    & 10 \\
  Duración por sprint  & Entre 1 y 4 semanas \\
  Capacidad total      & 116 puntos \\
  \bottomrule
\end{tabularx}

% ---------- 2. Objetivo del proyecto ----------
\section{Objetivo del proyecto}

\begin{objetivo}
  Validar el problema, definir el modelo de negocio, calcular el mercado y construir un MVP funcional de CrediStep, un GPS financiero que acompaña a los jóvenes colombianos antes de su primer crédito, y sustentar el proyecto en un pitch final ante los jurados.
  
\end{objetivo}

% ---------- 3. Backlog del producto ----------
\section{Backlog del producto}

\subsection{Épicas del proyecto}

\begin{tabularx}{\textwidth}{c X c c c}
  \toprule
  \textbf{ID} & \textbf{Épica} & \textbf{Sprint} & \textbf{Puntos} & \textbf{Estado} \\
  \midrule
  EP-01 & Como Product Owner, quiero el problema documentado con fuentes oficiales, para saber a quién ayuda CrediStep. & 1 & 10 & Terminada \\
  EP-02 & Como jurado, quiero un primer pitch de la idea, para entender el problema y la solución. & 2 & 10 & Terminada \\
  EP-03 & Como Product Owner, quiero el problema validado con una encuesta a jóvenes, para decidir si el proyecto sigue. & 3 & 13 & Terminada \\
  EP-04 & Como evaluador del proyecto, quiero el modelo de negocio descrito en un lienzo Canvas, para entender cómo se sostiene el proyecto. & 4 & 8 & Terminada \\
  EP-05 & Como evaluador del proyecto, quiero el TAM, el SAM y el SOM calculados, para conocer el tamaño del mercado. & 5 & 13 & Terminada \\
  EP-06 & Como Product Owner, quiero el flujo principal de la aplicación bosquejado a mano, para ver cómo lo usaría un joven. & 6 & 8 & Terminada \\
  EP-07 & Como Product Owner, quiero el contexto financiero elaborado, para ver los ingresos y costos proyectados. & 7 & 10 & Terminada \\
  EP-08 & Como jurado, quiero un segundo pitch de tres minutos, para conocer la validación y el mercado. & 8 & 10 & Terminada \\
  EP-09 & Como Product Owner, quiero un MVP funcional y documentado, para ver la solución en funcionamiento. & 9 & 18 & Terminada \\
  EP-10 & Como Product Owner, quiero el pitch final ante jurados y los entregables enviados, para cerrar el proyecto. & 10 & 16 & En curso \\
  \bottomrule
\end{tabularx}

% ---------- 4. Plan de sprints ----------
\section{Plan de sprints}

\begin{tabularx}{\textwidth}{c X l l c}
  \toprule
  \textbf{Sprint} & \textbf{Tema} & \textbf{Inicio} & \textbf{Fin} & \textbf{Historias} \\
  \midrule
  1 & Contextualización del problema & 16/07/2026 & 30/07/2026 & 3 \\
  2 & Primer pitch & 03/08/2026 & 31/08/2026 & 3 \\
  3 & Artefactos de validación & 01/09/2026 & 21/09/2026 & 3 \\
  4 & Lienzo Canvas & 22/09/2026 & 05/10/2026 & 2 \\
  5 & Definición de mercado & 06/10/2026 & 12/10/2026 & 4 \\
  6 & Prototipado a mano & 13/10/2026 & 19/10/2026 & 3 \\
  7 & Contexto financiero & 20/10/2026 & 26/10/2026 & 3 \\
  8 & Segundo pitch & 27/10/2026 & 02/11/2026 & 3 \\
  9 & MVP y mockups & 03/11/2026 & 16/11/2026 & 5 \\
  10 & Pitch final & 17/11/2026 & 25/11/2026 & 4 \\
  \bottomrule
\end{tabularx}

\subsection{Entregables por sprint}

\begin{tabularx}{\textwidth}{c X l}
  \toprule
  \textbf{Sprint} & \textbf{Entregable principal} & \textbf{Responsable} \\
  \midrule
  1 & Documento de contextualización con cifras y fuentes & Sarahy Silva \\
  2 & Pitch de nueve diapositivas con su guion & Valentina Carrillo \\
  3 & Encuesta a 126 jóvenes y análisis de resultados & Miguel Gómez \\
  4 & Lienzo Canvas con la explicación de sus bloques & Beiker Jaramillo \\
  5 & Definición de mercado, versión 1 & Laura Ortiz \\
  6 & Prototipo a mano de cinco pantallas & Sarahy Silva \\
  7 & Hoja de simulación financiera & Valentina Carrillo \\
  8 & Guion, diapositivas y video del segundo pitch & Miguel Gómez \\
  9 & Aplicación web, su documentación y el mercado en versión 2 & Beiker Jaramillo \\
  10 & Diapositivas del pitch final y entregables & Laura Ortiz \\
  \bottomrule
\end{tabularx}

% ---------- 5. Seguimiento del proyecto ----------
\section{Seguimiento del proyecto}

\begin{tabularx}{\textwidth}{c X X}
  \toprule
  \textbf{Sprint} & \textbf{Lo que se logró} & \textbf{Lo que quedó para después} \\
  \midrule
  1 & Problema y usuario definidos con datos oficiales & Guardar los PDF de las fuentes en la carpeta \\
  2 & Pitch que cubre las secciones esperadas & Manejar el tiempo y enviar el deck \\
  3 & Supuesto inicial validado con 126 respuestas & Enviar el Excel/CSV y las fechas de aplicación \\
  4 & Modelo de negocio completo en el lienzo & Reducir el texto del lienzo y corregir canales \\
  5 & Mercado calculado en su versión 1 & Ampliarlo más allá de Bogotá \\
  6 & Flujo de la aplicación en cinco pantallas & Consolidar los prototipos en un MVP \\
  7 & Base financiera con los precios de la aplicación & Alinear la hoja con el mercado \\
  8 & Pitch de tres minutos, en clase y en video & Incluir los aspectos financieros \\
  9 & MVP web funcional y mercado en versión 2 & Probarlo con usuarios reales \\
  10 & Primer borrador de las diapositivas & Ensayo, hoja financiera y entregables \\
  \bottomrule
\end{tabularx}

\subsection{Impedimentos y acciones}

\begin{tabularx}{\textwidth}{X l l}
  \toprule
  \textbf{Impedimento} & \textbf{Responsable} & \textbf{Estado} \\
  \midrule
  Rango de edad mal citado en una cifra (Sprint 1) & Valentina Carrillo & Resuelto \\
  Los títulos del guion y de las diapositivas no coincidían (Sprint 2) & Beiker Jaramillo & Resuelto \\
  El deck usado en el pitch presencial no se envió al profesor (Sprint 2) & Beiker Jaramillo & Abierto \\
  La muestra se concentró en jóvenes de 18 a 20 años (Sprint 3) & Miguel Gómez & Abierto \\
  No se envió el Excel/CSV exportado con el detalle de las respuestas (Sprint 3) & Miguel Gómez & Abierto \\
  La estructura de costos no tiene cifras (Sprint 4) & Laura Ortiz & Abierto \\
  La conversión a pago del 20 \% es un supuesto sin fuente (Sprint 5) & Laura Ortiz & Abierto \\
  La plantilla trae precios y productos que la aplicación no tiene (Sprint 7) & Laura Ortiz & Resuelto \\
  El pitch no incluyó los aspectos financieros (Sprint 8) & Valentina Carrillo & Abierto \\
  La hoja financiera no coincide con el mercado de la versión 2 (Sprint 9) & Laura Ortiz & Abierto \\
  La hoja financiera no coincide con las cifras de las diapositivas (Sprint 10) & Sarahy Silva & Abierto \\
  Falta el Excel/CSV con las respuestas de la encuesta (Sprint 10) & Miguel Gómez & Abierto \\
  Faltan las fotos del equipo en las diapositivas (Sprint 10) & Laura Ortiz & Abierto \\
  \bottomrule
\end{tabularx}

% ---------- 6. Métricas del proyecto ----------
\section{Métricas del proyecto}

\begin{tabularx}{\textwidth}{X c c}
  \toprule
  \textbf{Indicador} & \textbf{Planeado} & \textbf{Real} \\
  \midrule
  Sprints realizados              & 10 & 9 cerrados y 1 en curso \\
  Puntos de historia completados  & 116 & 100 \\
  Historias completadas           & 33 & 29 \\
  Velocidad promedio por sprint   & 11,6 & 11,1 \\
  \bottomrule
\end{tabularx}

\subsection{Puntos por sprint}

{\setlength{\tabcolsep}{4pt}
\begin{tabularx}{\textwidth}{X c c c c c c c c c c c}
  \toprule
  \textbf{Sprint} & \textbf{1} & \textbf{2} & \textbf{3} & \textbf{4} & \textbf{5} & \textbf{6} & \textbf{7} & \textbf{8} & \textbf{9} & \textbf{10} & \textbf{Total} \\
  \midrule
  Puntos planeados    & 10 & 10 & 13 & 8 & 13 & 8 & 10 & 10 & 18 & 16 & 116 \\
  Puntos completados  & 10 & 10 & 13 & 8 & 13 & 8 & 10 & 10 & 18 & -- & 100 \\
  \bottomrule
\end{tabularx}}

{El Sprint 10 termina el 25/11/2026; sus puntos se registran al cierre.}

% ---------- 7. Revisión del proyecto ----------
\section{Revisión del proyecto}

\subsection{Incrementos entregados}
\begin{itemize}[leftmargin=1.5em]
  \item {Problema respaldado con datos oficiales de inclusión financiera.}
  \item {Validación con 126 jóvenes: el 54 \% se sentiría inseguro al pedir su primer crédito y el 35,7 \% ya evitó pedirlo.}
  \item {Modelo de negocio con dos fuentes de ingreso: suscripción y licencias para universidades.}
  \item {Mercado calculado: 11,6 millones de jóvenes de 18 a 30 años (TAM), 4,2 millones atendibles (SAM) y 439 suscriptores de pago en el primer año (SOM), unos 105 millones de pesos.}
  \item {Prototipo a mano, MVP web funcional y su documentación.}
  \item {Dos pitches presentados y el borrador de las diapositivas del pitch final.}
\end{itemize}

\subsection{Retroalimentación de los interesados}
{Comentarios del profesor, que actuó como Product Owner, en las entregas calificadas:}

\begin{tabularx}{\textwidth}{c l c X}
  \toprule
  \textbf{Sprint} & \textbf{Entrega} & \textbf{Nota} & \textbf{Comentario principal} \\
  \midrule
  2 & Primer pitch & 3,8 & Cubre las secciones esperadas. Usamos casi el doble del tiempo, faltó presentar al equipo y un buen cierre, y no enviamos el deck. \\
  3 & Artefactos de validación & 4,6 & Buenas preguntas y análisis completo. Faltaron las fechas de aplicación y el Excel/CSV con el detalle de las respuestas. \\
  4 & Lienzo Canvas & 4,5 & Modelo bien estructurado. El lienzo tiene mucho texto, y parte de canales va en relacionamiento con clientes. \\
  6 & Prototipado a mano & 5,0 & Buen trabajo en el prototipo entregado por Beiker; muestra funcionalidades base para el MVP final. \\
  8 & Segundo pitch & 4,3 & Buen pitch en general. Faltaron los aspectos financieros, pidió subir más el mercado y el video se pasó un minuto. \\
  \bottomrule
\end{tabularx}

{El mercado se recalculó en el Sprint 9 por ese comentario: el SOM pasó de una sola universidad a 439 suscriptores de pago a escala nacional. Los sprints 1, 5, 7 y 9 no tienen una nota registrada en este documento, y el pitch final se evalúa el 24/11/2026.}

\subsection{Pendientes del proyecto}
\begin{tabularx}{\textwidth}{c X X}
  \toprule
  \textbf{Sprint} & \textbf{Pendiente} & \textbf{Acción propuesta} \\
  \midrule
  2 & El deck del primer pitch no se envió & Incluirlo en los entregables finales \\
  3 & Faltan las fechas de aplicación y el Excel/CSV de respuestas & Enviarlos en la entrega final \\
  4 & El lienzo tiene mucho texto y mezcla canales con relacionamiento & Ajustar el lienzo con frases cortas \\
  7 & La hoja financiera no coincide con el mercado de la versión 2 & Actualizarla antes del pitch final \\
  7 & La licencia institucional no tiene precio & Definirlo con la universidad \\
  9 & El MVP no se ha probado con usuarios reales & Probarlo con cinco jóvenes \\
  10 & Pitch final y entregables & Presentar el 24/11 y enviar el 25/11 \\
  \bottomrule
\end{tabularx}

% ---------- 8. Retrospectiva ----------
\section{Retrospectiva del proyecto}

\subsection{Qué salió bien}
\begin{itemize}[leftmargin=1.5em]
  \item {Cada sprint dejó un entregable concreto, y los nueve primeros se completaron.}
  \item {El problema se sostuvo con datos oficiales y con una encuesta propia.}
  \item {La idea central se mantuvo igual de principio a fin: aprender antes de cometer el error.}
  \item {Los comentarios del profesor se convirtieron en cambios: mercado más amplio, MVP funcional y aspectos financieros en el pitch final.}
\end{itemize}

\subsection{Qué se puede mejorar}
\begin{itemize}[leftmargin=1.5em]
  \item {Varios documentos se escribieron después de usarlos; la contextualización, por ejemplo, llegó después del primer pitch.}
  \item {Cuando una cifra cambió, no siempre se actualizó en los demás documentos, como pasó con el mercado y la hoja financiera.}
  \item {El manejo del tiempo falló en los dos pitches: el primero duró casi el doble y el video del segundo se pasó un minuto.}
  \item {Faltaron evidencias y piezas que el profesor pidió: el deck, el archivo de respuestas y los aspectos financieros.}
\end{itemize}

\subsection{Acciones de mejora}
\begin{tabularx}{\textwidth}{X l l}
  \toprule
  \textbf{Acción} & \textbf{Responsable} & \textbf{Fecha límite} \\
  \midrule
  Presentar el pitch final ante los jurados & Beiker Jaramillo & 24/11/2026 \\
  Actualizar la hoja de proyección financiera & Sarahy Silva & 24/11/2026 \\
  Ensayar el pitch con cronómetro y con la demo & Valentina Carrillo & 24/11/2026 \\
  Probar el MVP con cinco jóvenes & Laura Ortiz & 24/11/2026 \\
  Enviar el Excel/CSV de respuestas y las fechas de aplicación & Miguel Gómez & 25/11/2026 \\
  Ajustar el lienzo: frases cortas y canales corregidos & Beiker Jaramillo & 25/11/2026 \\
  Enviar los entregables finales por Bright Space & Miguel Gómez & 25/11/2026 \\
  \bottomrule
\end{tabularx}

% ---------- 9. Aprobaciones ----------
\section{Aprobaciones}

\vspace{1cm}
\begin{tabularx}{\textwidth}{X X}
  \textbf{Product Owner} & \textbf{Scrum Master} \\[2cm]
  \rule{6cm}{0.4pt} & \rule{6cm}{0.4pt} \\
  Firma y fecha & Firma y fecha \\
\end{tabularx}

\end{document}
